% Part: incompleteness
% Chapter: incompleteness-provability
% Section: provability-conditions

\documentclass[../../../include/open-logic-section]{subfiles}

\begin{document}

\olfileid[fr]{inc}{inp}{prc}

\olsection{Les conditions de dérivabilité pour $\Th{PA}$}

L'arithmétique de Peano, $\Th{PA}$, étend
$\Th{Q}$ en lui ajoutant les axiomes d'induction
pour toutes les !!{formula}s. Autrement dit, pour
chaque !!{formula}~$!A$, nous ajoutons la clôture universelle sur ses paramètres libres
d'une instance de la forme
\[
(!A(0)\land\lforall[x][(!A(x)\lif !A(x'))])
\lif\lforall[x][!A(x)].
\]
Il s'agit bien d'un \emph{schéma}, donc d'une
infinité d'axiomes; $\Th{PA}$ n'est d'ailleurs
pas finiment axiomatisable. Toutefois, on peut
décider effectivement si une chaîne de symboles
est une instance du schéma d'induction. L'ensemble
des axiomes de~$\Th{PA}$ est donc calculable.
Cette théorie est beaucoup plus puissante que
$\Th{Q}$. Par exemple, l'induction usuelle
permet d'y démontrer facilement la commutativité
de l'addition et de la multiplication. La plupart
des raisonnements finitaires d'arithmétique et de
combinatoire peuvent même s'y formaliser.

Puisque $\Th{PA}$ admet une axiomatisation
calculable, la relation de vérification des preuves
$\Prf[\Th{PA}](x,y)$ est calculable et donc
représentable dans~$\Th{Q}$, et a fortiori dans
$\Th{PA}$. Comme auparavant,
$\OPrf[\Th{PA}](x,y)$ désignera la !!{formula}
qui représente cette relation. Posons
$\OProv[\Th{PA}](y)=
\lexists[x][\OPrf[\Th{PA}](x,y)]$.
Intuitivement, cette !!{formula} dit que la !!{formula} de code~$y$ est démontrable à partir
des axiomes de~$\Th{PA}$. Nous avons besoin d'une
théorie un peu plus forte que~$\Th{Q}$ afin de
démontrer certains faits fondamentaux sur ce
prédicat de démontrabilité. Pour des !!{formula}s $!A$ et $!B$, voici
les conditions qu'il nous faut :
\begin{enumerate}
\item[P1.] Si $\Th{PA}\Proves !A$, alors
  $\Th{PA}\Proves\OProv[\Th{PA}](\gn{!A})$.
\item[P2.] Pour toutes !!{formula}s $!A$ et $!B$,
  \[
  \Th{PA}\Proves\OProv[\Th{PA}](\gn{!A\lif !B})
  \lif(\OProv[\Th{PA}](\gn{!A})
  \lif\OProv[\Th{PA}](\gn{!B})).
  \]
\item[P3.] Pour toute !!{formula}~$!A$,
  \[
  \Th{PA}\Proves\OProv[\Th{PA}](\gn{!A})
  \lif\OProv[\Th{PA}]
  (\gn{\OProv[\Th{PA}](\gn{!A})}).
  \]
\end{enumerate}
Pour P1, une démonstration effective de~$!A$ fournit un code
numérique précis. La représentabilité donne alors une preuve
de l'instance de~$\OPrf[\Th{PA}]$ portant sur ce code et celui
de~$!A$. L'introduction du quantificateur existentiel donne
$\OProv[\Th{PA}](\gn{!A})$. Pour P2, il faut formaliser
dans~$\Th{PA}$ l'opération qui combine une démonstration
de~$!A\lif !B$ et une démonstration de~$!A$, puis ajoute
une application du modus ponens.

Pour P3, on choisit un prédicat usuel de démontrabilité de
classe~$\Sigma_1$. On utilise alors la complétude
$\Sigma_1$ formalisée dans~$\Th{PA}$ : pour chaque énoncé
$!C$ de cette classe, $\Th{PA}$ démontre
$!C\lif\OProv[\Th{PA}](\gn{!C})$. On l'applique à
$!C\ident\OProv[\Th{PA}](\gn{!A})$. Une représentation
ponctuelle de la relation de preuve suffit pour l'argument de
P1, mais P2 et P3 demandent des vérifications internes à
$\Th{PA}$. Nous admettrons ici ces vérifications pour le
codage usuel des démonstrations.
Un choix raisonnable de $\OProv[\Th{PA}](y)$
satisfait également :
\begin{enumerate}
\item[P4.] Si
  $\Th{PA}\Proves\OProv[\Th{PA}](\gn{!A})$,
  alors $\Th{PA}\Proves !A$.
\end{enumerate}
Pour le prédicat usuel, la correction de~$\Th{PA}$ dans le
modèle standard justifie P4. Si la théorie démontre l'énoncé
$\OProv[\Th{PA}](\gn{!A})$, il existe réellement un code
de démonstration de~$!A$. P4 est une règle au sujet de ce que
la théorie démontre. Les deux théorèmes qui suivent utilisent
seulement P1--P3.

\begin{digress}
Gödel lui-même a été, sur ce point, aussi peu
disert que nous. À la fin de son article de 1931,
il esquisse la preuve du second théorème
d'incomplétude et promet d'en donner les détails
dans un article ultérieur. Il ne l'a jamais fait :
les lecteurs qui comprenaient l'argument étaient
convaincus qu'il pouvait être mené à bien, si bien
qu'il n'avait pas besoin d'en rédiger les détails.
\end{digress}

\end{document}
